% Context from chapters/wavelets.tex; for mathematical reference, not compilation.
}{wavcoef-1d-2}&
\image{wavelets}{.48}{wavcoef-1d-3}\\
$d_{-6}$ & $d_{-5}$ 
}
}{ 
	Wavelet coefficients. Top: the complete decomposition. Below: enlarged views of individual scales. 
}{fig-wavcoef-1d}


   
\subsection{Forward Fast Wavelet Transform (FWT)}

The algorithm applies a sequence of elementary operators
\eql{\label{eq-wavelet-step}
	\foralls j = J+1, \ldots, 0, \quad 
	(a_{j},d_{j}) = \Ww_j(a_{j-1})
	\qwhereq
	\Ww_j : \RR^{2^{-j+1}} \rightarrow \RR^{2^{-j}} \times \RR^{2^{-j}}
}
There are $|J|=\log_2(N)$ steps. Each $\Ww_j$ changes coordinates between orthonormal bases and is therefore orthogonal.

To describe the algorithm for $\Ww_j$, we introduce the filter coefficients $h,g \in \RR^\ZZ$
\eql{\label{eq-dfn-h-g}
	h_n \eqdef \frac{1}{\sqrt{2}} \dotp{ \phi(\cdot/2) }{\phi(\cdot-n)}
	\qandq
	g_n \eqdef \frac{1}{\sqrt{2}} \dotp{ \psi(\cdot/2) }{\phi(\cdot-n)}.
}
Figure~\ref{fig-haar-inner} illustrates the computation of these weights for the Haar system.

\begin{figure}[tbp]
\centering
\BookDrawing{tikz/wavelets-haar-inner-products}
\caption{Haar filter weights as inner products.}
\label{fig-haar-inner}
\end{figure}


We denote by $\downarrow_2 : \RR^K \rightarrow \RR^{K/2}$ the subsampling operator by a factor of 2, i.e.
\eq{
	u\downarrow_2 \eqdef (u_0,u_2,u_4,\ldots,u_{K-2}).
}
Throughout the following derivation, the scaling function, wavelet, and filters are real, and the filters decay sufficiently fast. The notation $\bar h_n=h_{-n}$ and $\bar g_n=g_{-n}$ denotes reflection. Complex filters would require conjugation as well.

\begin{prop}
One has
\eql{\label{eq-convol-subsample}
	\Ww_j(a_{j-1})= ( (\bar h \star a_{j-1})\downarrow_2, (\bar g \star a_{j-1})\downarrow_2 ),
}
where $\star$ denotes periodic convolutions on $\RR^{2^{-j+1}}$.
\end{prop}

\begin{proof}
We note that $\phi(\cdot/2)$ and $\psi(\cdot/2)$ belong to $V_0$, so one has the decompositions
\eql{\label{eq-phipsi-recurse-1}
	\frac{1}{\sqrt{2}} \phi(t/2) = \sum_n h_n \phi(t-n)
	\qandq
	\frac{1}{\sqrt{2}} \psi(t/2) = \sum_n g_n \phi(t-n)
}
Making the substitution $t \mapsto \frac{t-2^j p}{2^{j-1}}$ in~\eqref{eq-phipsi-recurse-1}, one obtains
\eq{
	\frac{1}{\sqrt{2}} \phi\pa{ \frac{t-2^j p}{2^j} }  = \sum_n h_n \phi\pa{ \frac{t}{2^{j-1}}-(n+2p)}
}
(similarly for $\psi$) and then reindexing $n \mapsto n-2p$, one obtains
\eq{
	\phi_{j,p} = \sum_{n \in \ZZ} h_{n-2p} \phi_{j-1,n}
	\qandq 
	\psi_{j,p}=\sum_{n\in\ZZ}g_{n-2p}\phi_{j-1,n}.
}
For periodized functions $(\phi_{j,n}^P,\psi_{j,p}^P)$, these identities remain valid after periodizing the filters modulo $2^{-j+1}$.
 
Taking inner products with $f$ on both sides gives the fundamental recursion below; the decay of $h,g$ justifies exchanging the sum and inner product.
\eql{\label{eq-recurse-formula-wav}
	a_{j,p} = \sum_{n \in \ZZ} h_{n-2p} a_{j-1,n} = (\bar h \star a_{j-1})_{2p}
	\qandq 
	d_{j,p} = \sum_{n \in \ZZ} g_{n-2p} a_{j-1,n} = (\bar g \star a_{j-1})_{2p}
}
where $\bar u_n \eqdef u_{-n}$. 
 
On the periodic domain $\TT=\RR/\ZZ$, the same recursion holds with $\star$ interpreted as convolution on $\ZZ/2^{-j+1}\ZZ$.
\end{proof}

Figure~\ref{fig-filterbank-fwd} shows two successive decomposition steps, each separating a coarse approximation from its details.

\myfigure{
\image{wavelets}{.95}{filterbank-fwd}
}{ 
	Forward filter bank decomposition.  
}{fig-filterbank-fwd}

The FWT thus operates as follows. The same real filters also act on complex-valued input signals:
\begin{rs}
	\item \textbf{Input:} signal $f \in \CC^N$.
	\item \textbf{Initialization:} $a_J = f$.
	\item \textbf{For} $j=J,\ldots,j_0-1$.
		\eq{
			a_{j+1} = (a_j \star \bar h) \downarrow 2 \qandq
			d_{j+1} = (a_j \star \bar g) \downarrow 2
		}
	\item \textbf{Output:} the coefficients $\{d_j\}_{J<j\leq j_0} \cup \{ a_{j_0} \}$.
\end{rs}

If the supports of $h$ and $g$ each contain at most $C$ indices, then evaluating each $\Ww_j$ requires $(2C) 2^{-j}$ operations, so that the complexity of the whole wavelet transform is
\eq{
	\sum_{j=J+1}^{0}(2C)2^{-j}=2C(2^{-J}-1)<2CN.
}
This shows that the fast wavelet transform is a linear-time algorithm.
 
Figure~\ref{fig-filterbank-pyramidal} shows the iterative extraction of the wavelet coefficients.
 
Figure~\ref{fig-wavalgo-1d} illustrates the storage layout: at each step, the new vectors $a_j$ and $d_j$ occupy the left portion of the output array.

\myfigure{
\image{wavelets}{.4}{filterbank-pyramidal}
}{ 
	Pyramidal computation of wavelet coefficients.  
}{fig-filterbank-pyramidal}

\myfigure{
\tabdeux{
\image{wavelets}{.48}{wavalgo-1d-0}&
\image{wavelets}{.48}{wavalgo-1d-1}\\
\image{wavelets}{.48}{wavalgo-1d-2}&
\image{wavelets}{.48}{wavalgo-1d-3}
}
}{ 
	Wavelet decomposition algorithm.  
}{fig-wavalgo-1d}


  
\paragraph{Fast Haar transform.}

For the Haar wavelets, one has
\eq{
	\phi_{j,n} = \frac{1}{\sqrt{2}} ( \phi_{j-1,2n} + \phi_{j-1,2n+1} ),
}
\eq{
	\psi_{j,n} = \frac{1}{\sqrt{2}} ( \phi_{j-1,2n} - \phi_{j-1,2n+1} ).
}
This corresponds to the filters
\eq{
	h = [\ldots,\;0,\;h[0]=\frac{1}{\sqrt{2}},\;\frac{1}{\sqrt{2}},\;0,\ldots],
}
\eq{
	g = [\ldots,\;0,\;g[0]=\frac{1}{\sqrt{2}},\;-\frac{1}{\sqrt{2}},\;0,\ldots].
}
The Haar transform iterates normalized sums and differences:
\begin{rs}
	\item \textbf{Input:} signal $f \in \CC^N$.
	\item \textbf{Initialization:} $a_J = f$.
	\item \textbf{For} $j=J,\ldots,j_0-1$.
		\eq{
			a_{j+1,n} = \frac{1}{\sqrt{2}} (a_{j,2n}+a_{j,2n+1}) \qandq
			d_{j+1,n} = \frac{1}{\sqrt{2}} (a_{j,2n}-a_{j,2n+1}).
		}
	\item \textbf{Output:} the coefficients $\{d_j\}_{J<j\leq j_0} \cup \{ a_{j_0} \}$.
\end{rs}

                                                            
\subsection{Inverse Fast Transform (iFWT)}

The inverse algorithm proceeds by inverting each step~\eqref{eq-wavelet-step}
\eql{\label{eq-wavelet-inverse-step}
	\foralls j=0,-1,\ldots,J+1, \quad
	a_{j-1} = \Ww_j^{-1}(a_{j},d_{j}) = \Ww_j^*(a_{j},d_{j}), 
}
where $\Ww_j^{*}$ is the adjoint for the canonical inner product on $\RR^{2^{-j+1}}$; in matrix form, it is the transpose.

Let $\uparrow_2 : \RR^{K/2} \rightarrow \RR^K$ denote up-sampling by inserting zeros:
\eq{
	a\uparrow_2 = (a_0,0,a_1,0,\ldots,a_{K/2-1},0) \in \RR^K.
}

\begin{prop}
One has
\eq{
	\Ww_j^{-1}(a_{j},d_{j}) = (a_{j}\uparrow_2) \star h +  (d_{j}\uparrow_2) \star g.
}
\end{prop}

\begin{proof}
Since $\Ww_j$ is orthogonal, $\Ww_j^{-1} = \Ww_j^*$.
 
We write the whole transform as 
\eq{
	\Ww_j = S_2 \circ C_{\bar h,\bar g} \circ \Dd
	\qwhereq
	\choice{
		\Dd(a)=(a,a), \\
		C_{\bar h,\bar g}(a,b)=(\bar h\star a,\bar g\star b), \\
		S_2( a,b ) = (a\downarrow_2, b\downarrow_2).	
	}
}
The adjoints are
\eq{
	\Dd^*(a,b)=a+b, \quad
	C_{\bar h,\bar g}^*(a,b) = (h \star a, g \star b), \qandq
	S_2^*(a,b)=(a\uparrow